\chapter{总结与展望}
\label{sec:conclusion}

\section{作品完成情况}

OralSage 面向口腔医学生对专业知识学习、专题复习和日常计划管理的实际需求，构建了由学习画像、主题学习、课程规划、资源生成、智能讲解、即时答疑、学习测评和计划执行组成的连续学习闭环。学生可以从学习主页面输入课程主题、生成学习计划或检索专业文献，在课程学习过程中查看讲义和多模态资源、进行随堂练习并随时提出问题；课程进度、提问内容和答题结果继续进入个人学习状态，为后续复习和计划调整提供依据。ClawBot 将网页端形成的学习计划延伸到微信聊天环境，使今日计划查看、任务打卡、总结回顾和任务调整能够在日常使用场景中继续完成。

在技术实现上，作品已经形成口腔医学 KAG 知识底座、多智能体协同链路、学习者状态模型和同主题资源生成体系。Planner、Resource Factory、Tutor、QA、Assessment 和 Memory/User Model 在 Orchestrator 的调度下分别承担课程规划、讲义生成、课程讲解、即时答疑、学习测评和状态更新。课程、问答和练习共享知识节点与资料片段，模型生成结果需要经过证据标识和数据结构校验，从而使专业知识检索、生成内容和学习状态保持关联。

当前知识底座覆盖 14 门口腔学科，包含 1,390 个核心条目、9,907 个知识概念、5,758 条知识关系和 3,161 道关联题目；统一知识资源中包含 21,737 个口腔医学文本单元和 34,805 条检索语料。围绕同一课程节点，系统能够生成讲义、Tutor 讲解、思维导图、PPT/slides、音频、练习题和阶段卷，并将测评结果回写到学习状态。

作品已完成自动化测试和真实功能链路测试。Python 后端测试共执行 961 项且无失败项；学习计划草稿、主对话计划草稿和主对话测试分别完成 60、76 和 51 项，学习计划逻辑测试完成全部测试项。以“牙周炎”为主题的真实流程生成 5 个课程节点，首节课程节点绑定 1 个知识节点和 2 条资料片段，首节讲义包含 177 字概述和 4 个主要知识点。当前章节答疑使用 2 条证据直接回答，课程延伸答疑生成 2 条定向检索计划并使用 4 条证据完成回答。上述结果表明，OralSage 已完成从专业知识检索到课程学习、练习反馈和计划执行的主要产品与技术链路，具备本地运行、功能演示和作品交付条件。

\section{当前不足}

OralSage 已经完成主要功能闭环，但口腔医学知识、生成式学习资源和长期个性化服务仍具有持续演进空间。现阶段的边界主要体现在以下三个方面。

\begin{enumerate}
  \item \textbf{专业知识覆盖需要持续扩展。}

  当前知识底座已经覆盖主要口腔学科，并形成教材、知识节点、知识关系、题目和检索语料的统一组织方式。随着教材版本更新、临床指南修订和前沿研究积累，系统仍需继续补充课程章节和专业资料，并维护不同来源内容之间的版本关系。现有数据处理和索引流程已经为后续扩展提供基础。

  \item \textbf{复杂资源的首次生成效率仍有优化空间。}

  课程规划、讲义生成和 Tutor 讲解需要依次完成知识检索、证据绑定、内容生成和结构校验，因此完整首节学习流程仍需要一定等待时间。在“牙周炎”真实链路测试中，课程大纲于 145.088 秒就绪，Tutor 首次内容于 195.103 秒就绪。当前系统已经采用进度事件、讲义缓存、同步生成和后台任务等机制，后续仍可通过并行生成、缓存复用和任务调度进一步缩短首次等待时间。

  \item \textbf{长期学习效果需要在持续使用中进一步积累。}

  当前测试已经验证学习画像、答题记录、掌握度和薄弱点能够随学习行为发生变化。例如，测试画像的薄弱点由空集逐步更新为“机制链条不完整”“牙槽骨吸收”和“牙龈炎与牙周炎区别”，讲解偏好和学习节奏也能够随过程信息调整。随着实际使用周期增长，系统还可以积累更长时间范围内的计划执行、知识保持和复习表现数据，以进一步评估个性化策略对学习效率的持续影响。
\end{enumerate}

\section{后续规划}

后续工作将在现有学习闭环基础上，从专业知识、生成效率、个性化能力、教师支持和跨端应用五个方向继续完善。

\begin{enumerate}
  \item \textbf{持续扩展专业知识空间。}

  增加更多口腔医学课程章节、临床指南和前沿文献，完善知识来源、更新时间和版本关系的管理，使新增资料能够继续进入课程规划、问答和练习链路。

  \item \textbf{优化多模态资源生成效率。}

  进一步完善课程节点级缓存、并行资源生成和任务优先级调度，优先返回当前学习所需的首节讲义和讲解，再在后台生成后续课件、音频和练习资源，缩短学生开始学习前的等待时间。

  \item \textbf{深化学习状态与个性化推荐。}

  利用长期学习记录改进掌握度估计、薄弱点识别、复习间隔和任务推荐，使课程讲解方式、练习难度和计划安排能够更加稳定地反映学习者的实际变化。

  \item \textbf{完善教师支持能力。}

  在现有课程资源和学习状态基础上，增加资源审核、学习进度查看、薄弱知识汇总和形成性评价入口，帮助教师了解学生在不同课程节点上的学习表现，并对生成内容和学习计划进行必要调整。

  \item \textbf{推进跨端学习场景应用。}

  继续完善网页端与 ClawBot 的任务同步和状态回写，使课程学习、任务提醒、计划打卡和学习总结在不同终端之间保持一致，并在更多真实学习场景中验证长期使用体验。
\end{enumerate}

OralSage 以口腔医学专业知识为基础，将课程规划、资源生成、智能讲解、即时答疑、学习测评和计划执行连接为连续学习过程。作品已经完成主要功能与技术链路，并通过自动化测试、真实检索测试和核心学习流程测试验证了系统可用性。后续工作将在现有闭环基础上继续扩展知识覆盖、提升生成效率并积累长期学习数据，使系统逐步成为能够持续支持口腔医学生学习、复习和自我管理的智能学习伙伴。
